The cascade exposes explicit trade\hyp offs: lowering $\tau\_{low}$ and/or raising $\tau\_{high}$ reduces false negatives but increases Stage~2 usage; conversely, tighter thresholds reduce escalation at the risk of missed fakes. Our tuning results indicate that, on the combined validation set, we can achieve sub\hyp percent FNR at Stage~2 escalation below 2\% (Table~\ref{tab:deployment_scenarios}), while cross\hyp dataset results on Deepfake\hyp Eval\hyp 2024 reveal distribution shift and robustness gaps that merit further study.

From an operational perspective, the key benefit of MobileDeepfake is not only higher AUC but also the availability of explicit dials: Stage~1 leakage $\pi\_{\mathrm{leak}}$, Stage~2 escalation $\pi\_2$, and calibrated thresholds/temperatures can be tuned to match risk budgets. For example, a platform prioritising safety may accept a modest increase in $\pi\_2$ to drive $\pi\_{\mathrm{leak}}$ and FNR down, whereas a bandwidth\hyp constrained device may choose a more conservative operating point and rely on human review for borderline cases. The ability to quantify these trade\hyp offs and to log them over time makes the system easier to audit and maintain than opaque single\hyp model detectors.

Compared to a single, stronger model deployed on all inputs, the cascade behaves like a compute\hyp aware wrapper that makes risk/efficiency trade\hyp offs explicit. A single EfficientNetV2 or ViT\hyp style detector~\cite{dosovitskiy2021vit,liu2021swin,pei2024deepfakesurvey} can achieve strong AUC but forces operators to choose a fixed operating point that simultaneously governs all samples, often leading to either high FNR (if thresholds are conservative) or high system load and false positives (if thresholds are aggressive). In our system, Stage~1 absorbs most easy real cases at very low cost, and Stage~2 is reserved for the ambiguous band; escalation rates on in\hyp distribution data are only a few percent, so the average per\hyp frame cost is dominated by the lightweight filter. On Deepfake\hyp Eval\hyp 2024, the Stage~2\hyp only configuration improves recall but at the price of extremely high FPR and nearly 100\% Stage~2 usage, whereas the calibrated cascade maintains bounded compute and a clearer understanding of where it fails.

\subsection*{Deployment Scenarios and Operating Points}
To make these trade\hyp offs more concrete, we outline representative operating points derived from Stage~4 grid search on the combined validation set (Table~\ref{tab:deployment_scenarios}). Each scenario corresponds to a pair of thresholds $(\tau\_{low},\tau\_{high})$ and reports the resulting false negative rate (FNR) and Stage~2 escalation rate $\pi\_2$; all other details (datasets, manifests, calibration) follow the protocol in Sections~\ref{tab:dataset_sizes} and~\ref{tab:cascade_best}.
\begin{table}[h]
  \centering
  \caption{Example deployment scenarios and operating points on the combined validation split (derived from Stage~4 best configurations). FNR and Stage~2 rate are reported as fractions.}
  \label{tab:deployment_scenarios}
  {\scriptsize
  \setlength{\tabcolsep}{4pt}
  \begin{tabular}{@{}lccccp{0.4\textwidth}@{}}
    \toprule
    Scenario & $\tau\_{low}$ & $\tau\_{high}$ & FNR & Stage~2 rate $\pi\_2$ & Notes \\
    \midrule
    Safety\hyp first moderation & 0.05 & 0.55 & 0.0060 & 0.0116 & Very low FNR, minimal expert usage \\
    Balanced mobile deployment & 0.03 & 0.55 & 0.0054 & 0.0150 & Near\hyp zero FNR with modest compute \\
    Device\hyp first offline tool & 0.10 & 0.57 & 0.0363 & 0.0859 & Lower compute, accepts higher FNR \\
    \bottomrule
  \end{tabular}}
\end{table}
In high\hyp security scenarios (e.g., platform\hyp side moderation pipelines), configurations like the ``safety\hyp first'' row keep FNR around 0.6\% while sending roughly 1\% of samples to the expert; this is suitable when false negatives are costly and infrastructure can absorb some Stage~2 load. For battery\hyp and bandwidth\hyp constrained offline tools, higher $\tau\_{low}$ and/or $\tau\_{high}$ reduce $\pi\_2$ at the cost of higher FNR; users can be warned when samples fall into an ``uncertain'' band and encouraged to seek additional verification. Hybrid settings such as fact\hyp checking workflows can adopt the balanced configuration and treat escalated cases as priority items for human review, using logged $\pi\_2$ and FNR over time as monitoring signals.

In our current implementation, dynamic/adaptive thresholds (per\hyp device or per\hyp family) are treated as future work rather than part of the main evaluated system. The public codebase includes a prototype hook for adjusting $(\tau\_{low},\tau\_{high})$ based on heuristic uncertainty scores, but all quantitative results reported in this paper use static thresholds selected by the Stage~4 grid search on cached logits.

\subsection*{PC vs.\ Mobile Deployment Consistency}
A key practical question is whether cascade behaviour observed on desktop transfers reliably to mobile deployment. Our on\hyp device evaluation on the \texttt{mobile\_eval} subset (Section~\ref{sec:ondevice_results}) demonstrates that using the same tuned threshold configuration from Stage~4 (e.g., $\tau\_{low}=0.05$, $\tau\_{high}=0.55$, $\tau\_{s2}=0.5$ for the safety\hyp first operating point) produces qualitatively similar results across platforms: both PC and mobile achieve low FNR ($<$3\%), confirming that the safety\hyp oriented bias is preserved through ONNX export.

The higher Stage~2 escalation rate on mobile (7.2\% vs.\ 1.2\% on combined validation) and modestly higher FPR (12.5\% vs.\ 3\%) primarily reflect differences in sample composition rather than platform\hyp specific degradation. The compact \texttt{mobile\_eval} subset, by design, samples diverse sources and includes challenging cases that are more likely to fall into the ambiguous band. When evaluating on identically distributed samples, we expect the escalation rate and error profiles to match more closely.

From a practical standpoint, the 180\,ms average latency observed on a Snapdragon~8~Gen~2 device is sufficient for interactive single\hyp image analysis (approximately 5--6 images per second). For video processing at 1~fps sampling, the cascade can process up to 5 frames per second, enabling real\hyp time feedback even without GPU acceleration. These numbers should be interpreted as a conservative, CPU\hyp only ``worst\hyp case'' baseline: they demonstrate that the system remains usable even when device NPUs/GPUs are unavailable or misconfigured. Devices with older SoCs or lower thermal budgets may require sparser sampling or staged batching, while future work enabling NNAPI or GPU delegates is expected to reduce latency further on capable hardware.

\subsection*{Application-Level Considerations}
The resulting FPR/FNR of the model on mobile (12.5\% FPR, 2.5\% FNR) means we flag one in eight real images as suspicious overall but miss about 1 in 40 fakes. If false positives can be reviewed by someone (e.g., in content moderation queues, or personal verification tools) then this tendency to err on the side of caution is acceptable - users are warned about possible manipulation, and a human can decide what to do with the borderline cases.

In use\hyp cases well\hyp suited to fully automated pipelines (e.g., blocking users at upload time), the model's 12.5\% FPR may be too high without further filtering (e.g., inferring user reputation scores or context signals). Since lowering FPR will increase FNR, the application context should dictate the ``sweet spot'' for when fraud detection flags should appear, and Blitz will allow you to choose your operating spot using the explicit threshold dials in the cascade without needing to retrain the model.

\subsection*{Future On-Device Work}
There are a few things we can do to enhance on\hyp device performance and broaden your use\hyp cases:
\begin{itemize}[leftmargin=*,itemsep=2pt]
  \item Multi\hyp device evaluation. Testing on mid\hyp range and entry\hyp level SoCs (e.g., Snapdragon~6/7 series, MediaTek Dimensity) would characterise performance variability and inform minimum hardware recommendations.
  \item Hardware acceleration. In ONNX Runtime, enabling an NPU or GPU delegate, or NNAPI would reduce inference latency for Stage~2.
  \item Energy and thermal profiling. Measuring battery drain and device temperature during prolonged batch processing would provide practical deployment guidance for continuous\hyp scanning applications.
  \item Advanced face detection. Upgrade from the deprecated Android \texttt{FaceDetector} to MediaPipe or an ONNX face model so we can more robustly preprocess images that have multiple faces or non\hyp frontal poses.
  \item Further model compression. For our Stage~2 distro, quantization\hyp aware training (QAT) or distillation to a smaller backbone would make our APK bag smaller and faster while keeping accuracy up.
\end{itemize}

Remaining limitations include residual domain shift, sensitivity to degradations, and dependence on calibration. Future directions include, (i) robustness stress testing and adaptive thresholds (per\hyp family, per\hyp device) learn about families of performance; (ii) distillation or architecture search learn about distilling or shrinking the amount of time the Stage~2 network takes to run through inference w/o compromising accuracy; and finally (iii) domain adaptation learn improving generalization w/o raising on\hyp device cost. Beyond technical metrics, longitudinal studies on real platforms - how do escalation rates change? How many more false negative incidents are seen? How does user experience vary? - would help concretely show that this approach does lead to better safety in deployed settings.

Our design trajectory originally considered a more aggressive compression path in which a large teacher (e.g., EfficientNetV2 or a hybrid CNN--Transformer) would be distilled into a single compact student model for on\hyp device deployment. Early prototypes confirmed that standard teacher--student distillation can preserve much of the teacher's AUC, but they also highlighted two limitations in our setting: first, a single distilled detector still exposes only one global operating point, making it difficult to express risk budgets and compute constraints explicitly; second, repeated distillation and re\hyp calibration cycles across datasets and perturbations added substantial engineering overhead. These observations motivated the current cascade: rather than pushing all complexity into a single highly optimised student, we keep a relatively simple Stage~1 filter, retain a moderately sized Stage~2 expert, and optimise thresholds and routing policies around them. Knowledge distillation thus becomes a complementary tool for future work on shrinking Stage~2, while the cascade provides the primary mechanism for exposing and tuning system\hyp level trade\hyp offs.

\subsection*{Domain Adaptation and Feature Diagnostics}
Feature\hyp space visualization (e.g., t\,\-SNE) and distributional distances (e.g., MMD) can quantify gaps between training sources and OOD sets~\cite{wang2018deepDA_survey}. Lightweight adaptation mechanisms---such as recalibrating temperatures and thresholds on small labelled OOD samples or mixing in high\hyp confidence pseudo\hyp labels---are promising deployment\hyp friendly options. In our preliminary experiments, however, naive pseudo\hyp labelling sometimes distorted calibration or overfit to a particular OOD snapshot. Designing a more principled adaptation protocol, for example via curriculum schedules that gradually mix real OOD labels with pseudo\hyp labels while explicitly monitoring FNR/FPR and Stage~2 usage per domain, remains an important direction for future work.

\subsection*{Early Exits and Overthinking}
Shallow\hyp Deep and related early\hyp exit works show that internal classifiers can both reduce compute and, in some cases, alter final predictions due to overthinking (early exits) \cite{kaya2019shallowdeep}. Our confidence\hyp based two\hyp threshold routing mitigates such risks by making Stage~1 decisions only under high confidence and escalating ambiguous cases, which we monitor via the Stage~2 escalation rate.

\subsection*{Threats to Validity}
\textbf{Data shifts.} Cross\hyp dataset gaps (e.g., Deepfake\hyp Eval\hyp 2024) may not fully reflect broader wild distributions. We mitigate via multi\hyp dataset training and difficult\hyp sample mining, but residual shift remains.

\textbf{Threshold sensitivity.} Results depend on grid\hyp searched low/high thresholds and (optionally) Stage~2 temperature. We report tuned values and provide robustness sweeps to reveal trade\hyp offs; deployment may require periodic re\hyp tuning.

\textbf{Perturbation scope.} Our suite covers JPEG, noise, blur, brightness/contrast, downscale, and gamma; other degradations (e.g., color casts, cropping, codecs) are not exhaustively tested.

\textbf{Compute constraints.} CPU\hyp only inference and sample\hyp limited sweeps are used for timely evaluation; exact numbers may vary under different hardware or larger samples. We prioritize trends and trade\hyp offs over absolute values.

\subsection*{GenConViT Integration: Lessons from a Hybrid CNN-Transformer Approach}

One of our exploratory directions involved integrating GenConViT~\cite{genconvit_github} as an alternative Stage~2 expert and as an additional feature source for the optional Stage~3 meta\hyp model. In our setting, the best GenConViT configurations matched but did not reliably exceed the EfficientNetV2\hyp B3 expert on the combined validation split, while requiring substantially longer training time and exhibiting higher variance across random seeds and hyperparameters. This combination of modest average gains, instability, and increased complexity made GenConViT a poor fit for our primary deployment target. As a result, we base the default cascade on a MobileNetV4 Stage~1 filter followed by an EfficientNetV2\hyp B3 Stage~2 expert, and treat GenConViT as a research component for future exploration. Additional details on the integration setup and experimental results are provided in the GenConViT appendix.

\subsection*{Other Partial Attempts and Lessons Learned}

\textbf{Meta\hyp features and fusion.} We prototyped several expert\hyp feature fusion strategies for the LightGBM meta\hyp model (e.g., single\hyp expert only, simple/weighted concatenation, normalization or variance\hyp guided selection). Early trials yielded small or inconsistent deltas over the strongest single expert; we thus report the most stable configuration and defer a full ablation to extended experiments.

\textbf{Routing alternatives.} Entropy/margin\hyp based routing were explored but not fully tuned under the same budgets; we keep them as a template for future runs and focus here on confidence\hyp based routing with explicit escalation control.

\textbf{Robustness trends.} Non\hyp monotonic patterns under JPEG/noise likely reflect sample size and threshold effects. We treat these as hypotheses rather than definitive claims, and plan larger sweeps with repeated runs and simple significance checks before drawing firm conclusions.
